% ============================================================
% DEEDS — SPECIALIZATIONS & SIGNATURE ABILITIES (first draft)
% ------------------------------------------------------------
% AUTHOR SPEC: 6 slots on the sheet; deeds are RANKED (I–IV),
% upgrading in place (Deadeye I -> II -> III -> IV). One deed
% occupies one slot at any rank.
%
% DESIGN GRAMMAR (consistent across all 20 deeds):
%   Rank I   — a specialization: +1 die or a small reliable edge
%              in the deed's exact territory.
%   Rank II  — a widening: a second edge, a floor, or a
%              once-per-scene trick.
%   Rank III — a signature: an active ability (Momentum-priced
%              or once per scene) that people describe when they
%              describe you.
%   Rank IV  — a legend: once per session (usually), the thing
%              the stories exaggerate — except they don't.
%
% [PROPOSALS FLAGGED FOR AUTHOR SIGN-OFF]
%   D1. Acquisition: one deed rank at every even character level
%       (new deed at I, or +1 rank to an owned deed), PLUS the GM
%       may award a rank when the table agrees a qualifying feat
%       occurred in play. Each deed lists "Earned by" flavor.
%   D2. Rank gates (placeholder): rank III requires level 5+;
%       rank IV requires level 9+.
%   D3. Stacking: deed dice stack with species traits, cards,
%       and augments; a character can't apply the same deed
%       twice to one roll.
%   D4. "Once per scene/session" language follows the core
%       chapter's definitions.
%   D5. Several rank-IV abilities are deliberately cinematic
%       (retroactive declarations, standoff interrupts); tune
%       per table.
%   D6. Iron in the Blood II references creature alignment tags
%       (Wild/Wyrd-aligned) — depends on the bestiary's tagging
%       scheme; placeholder until the Threats chapter.
% ============================================================

\chapter{Deeds}

\section{What a Deed Is}

The six lines under \textbf{Deeds} on your sheet are not a shopping list. They are the parts of your legend that have hardened into fact.

A Deed is a specialization or signature ability, earned in play and ranked \textbf{I through IV}. Each deed occupies one of your six slots at any rank --- you write \textit{Deadeye III} on the line, not three separate entries --- so a career holds at most six legends, and which six is as much a portrait as your Demeanor. Ranks are upgrades in place: the deed deepens; it does not multiply.

[PROPOSAL D1] \textbf{Earning deeds.} You gain one deed rank at every even character level --- a new deed at rank I, or one more rank in a deed you hold --- \textit{and} the GM may award a rank, at any time, when the table agrees the character just \textit{did the thing}: held the door, made the shot, talked the room around. Every deed below carries an ``Earned by'' line for exactly this purpose. The award track and the level track are the same track; deeds do not care how the legend started, only that it's true. [PROPOSAL D2] Rank III requires level 5+; rank IV requires level 9+.

[PROPOSAL D3] \textbf{Stacking.} Deed bonuses stack with species traits, cards, and Alignment augments. You cannot apply the same deed to a single roll twice.

\textbf{Reading the ranks.} Every deed follows the same grammar. \textbf{Rank I} is a specialization --- a die, an edge, a reliability, in the deed's exact territory. \textbf{Rank II} widens it. \textbf{Rank III} is a signature: the active ability people describe when they describe you. \textbf{Rank IV} is a legend --- usually once per session --- the thing the stories exaggerate, except they don't.

\medskip\hrule\medskip

\section{Deeds of Arms}

\subsection{Deadeye}
\textit{Earned by the shot nobody believed until they walked over and looked.}
\begin{description}
    \item[I.] Once per activation, reroll one missed die on a Shoot attack.
    \item[II.] Your Shoot attacks ignore 1 point of the target's Defend.
    \item[III.] \textbf{Dead to Rights} (1 Momentum, once per scene): declare before rolling --- this Shoot attack cannot be reduced below half your Shoot dice (round up) by Defend.
    \item[IV.] Once per session, after rolling a Shoot attack, every die counts as a hit.
\end{description}

\subsection{Duelist}
\textit{Earned by winning a fight you started by standing still.}
\begin{description}
    \item[I.] Once per activation, reroll one missed die on a Strike attack.
    \item[II.] Your Strike attacks ignore 1 point of the target's Defend.
    \item[III.] \textbf{Riposte} (free, once per round): when an adjacent enemy attacks you and scores no hits, immediately make a Strike attack against them at no Momentum cost.
    \item[IV.] Once per session, one Strike attack ignores Defend and Armor entirely.
\end{description}

\subsection{Bulwark}
\textit{Earned by being the reason the count of the dead stopped where it did.}
\begin{description}
    \item[I.] +1 Defend during any round in which you have not moved.
    \item[II.] Allies adjacent to you gain +1 Defend.
    \item[III.] \textbf{Interpose} (free, once per round): when an attack targets an adjacent ally, you become its target instead.
    \item[IV.] Once per session, reduce one attack against you or an adjacent ally to zero hits after the roll.
\end{description}

\subsection{Deathproof}
\textit{Earned by being carried off the field, and back onto the next one.}
\begin{description}
    \item[I.] +2 maximum HP.
    \item[II.] +2 more maximum HP, and while at half HP or below, gain +1 die on saves.
    \item[III.] Once per scene, when you would drop to 0 HP or below, drop to 1 instead.
    \item[IV.] Once per session, when you would die, you don't: you are at 1 HP, standing, and visibly annoyed.
\end{description}

\medskip\hrule\medskip

\section{Deeds of Tempo}

\subsection{Quickdraw}
\textit{Earned by clearing leather while the other party was still deciding.}
\begin{description}
    \item[I.] Once per scene, when an enemy targets you, play one 1-cost card from your hand as a reaction, at no Momentum cost, resolving it before their action.
    \item[II.] Twice per scene.
    \item[III.] Your reaction may instead be a Shoot attack, or a card costing up to 2.
    \item[IV.] \textbf{The Standoff} (once per session): when an enemy declares an attack against you, resolve a full Shoot attack against them first. If it drops them, their attack never happens.
\end{description}

\subsection{Pacesetter}
\textit{Earned by being, verifiably, the reason the whole side moved like that.}
\begin{description}
    \item[I.] Once per scene, add +1 to your side's Momentum roll.
    \item[II.] Your side's Momentum rolls of 1 count as 2.
    \item[III.] Your side gains +1 bonus Momentum on every surge you are conscious for.
    \item[IV.] Once per session, after your side rolls Momentum, set the die to 6.
\end{description}

\subsection{Ace in the Hole}
\textit{Earned by producing, at the fatal moment, exactly the card everyone forgot you had.}
\begin{description}
    \item[I.] When you load your hand at the start of an encounter, you may search your holstered pile for one card and load it as one of your five, then shuffle.
    \item[II.] Once per scene, retrieve one card (not a One-Shot) from your expended pile into your hand.
    \item[III.] Your loaded hand holds 6 cards.
    \item[IV.] Once per session, return one of your used One-Shot cards to your hand.
\end{description}

\subsection{Old Reliable}
\textit{Earned by doing the same beautiful thing so often it acquired your name.}
\begin{description}
    \item[I.] Choose one card in your deck: your \textit{reliable}. When you load your hand at the start of an encounter, you may take the reliable from holstered as one of your five.
    \item[II.] Once per scene, when the reliable enters your expended pile, return it to holstered immediately.
    \item[III.] After you play the reliable, immediately load one card from holstered (this can exceed your normal reload timing, not your hand size).
    \item[IV.] At the start of each of your side's surges, if the reliable is in your expended pile, return it to your hand.
\end{description}

\medskip\hrule\medskip

\section{Deeds of Wits}

\subsection{Sharp Eye}
\textit{Earned by seeing the thing whose whole job was not being seen.}
\begin{description}
    \item[I.] +1 die on Perceive.
    \item[II.] Once per scene, after examining a place, person, or object, ask the GM one direct question about it and receive a true answer.
    \item[III.] You notice as though searching: within 10 tiles, hidden things, ambushes, and wrongness prompt a Perceive roll from you even when you aren't looking.
    \item[IV.] \textbf{I Already Noticed} (once per session): retroactively declare that you clocked a detail earlier in the scene --- the exit, the second knife, the face in the crowd. The GM makes it true within reason.
\end{description}

\subsection{Silver Tongue}
\textit{Earned by walking out of a room that everyone agreed you would be carried out of.}
\begin{description}
    \item[I.] +1 die on Convince and Deceive.
    \item[II.] Once per scene, reroll a failed social roll by genuinely changing your approach.
    \item[III.] \textbf{Parley} (1 Momentum, once per scene): one sapient enemy who can hear you will not act against you or your allies until your side's next surge, or until any of you acts against them.
    \item[IV.] Once per session, one claim you make is believed by its audience until hard evidence contradicts it.
\end{description}

\subsection{Trailwise}
\textit{Earned by bringing everyone home through country that had other plans.}
\begin{description}
    \item[I.] +1 die on Survive.
    \item[II.] Companions traveling with you may use your Survive against travel hazards, weather, and forage.
    \item[III.] Sign speaks to you: automatically learn number, kind, gait, burden, and age from any trail you examine.
    \item[IV.] \textbf{There's Always a Way} (once per session): declare a plausible route feature --- a ford, a pass, a culvert under the line --- and the GM places it within reach.
\end{description}

\subsection{Learned}
\textit{Earned by knowing the answer in the field, under fire, without the library.}
\begin{description}
    \item[I.] +1 die on one of Study, Folklore, or Religion (chosen when this deed is earned).
    \item[II.] Once per scene, recall a relevant, useful, true fact from your discipline --- the GM supplies it.
    \item[III.] With time and materials, your research counts as two ranks higher, and takes half the time.
    \item[IV.] \textbf{The Briefing} (once per session): hold forth on a named threat or situation for a minute; your whole side gains +1 die on rolls against it for the rest of the scene.
\end{description}

\medskip\hrule\medskip

\section{Deeds of the Planes}

\subsection{Planewise}
\textit{Earned by standing in the current and not being carried.}
\begin{description}
    \item[I.] +1 die on all Alignment saves.
    \item[II.] Choose one plane: gain 1 automatic success on saves against it.
    \item[III.] Once per scene, reroll a failed Alignment save entirely.
    \item[IV.] Once per session, automatically succeed at one Alignment save, margin equal to your rating.
\end{description}

\subsection{Even Keel}
\textit{Earned by borrowing from all five natures and dying in debt to none.}
\begin{description}
    \item[I.] When you augment a Skill roll with 2 or fewer Alignment dice, the plane never complicates.
    \item[II.] +1 die whenever you augment (the plane leans in).
    \item[III.] Once per scene, augment a Skill roll \textit{after} seeing its result.
    \item[IV.] Once per session, augment one roll with two different planes at once, both at full rating.
\end{description}

\subsection{Deep Accounts}
\textit{Earned by paying the planes' prices in full, on time, with your eyes open.}
\begin{description}
    \item[I.] Once per scene, when you invoke an Alignment, gain +1 automatic success.
    \item[II.] Once per scene, when you invoke, defer the price: the complication arrives later this session, at a moment you choose (its content remains the GM's).
    \item[III.] \textbf{Name the Terms} (once per session): when you invoke, the GM offers two different prices; you choose which you pay.
    \item[IV.] Choose one plane: your rating counts as 2 higher (max 10) when invoking it. Its prices scale accordingly. The plane knows your name.
\end{description}

\subsection{Iron in the Blood}
\textit{Earned by walking out of a thin place that kept everyone else.}
\begin{description}
    \item[I.] +1 die on saves against Wild and Wyrd effects.
    \item[II.] [PROPOSAL D6] Your attacks count as cold iron: against Wild- and Wyrd-aligned creatures, ignore 1 Armor.
    \item[III.] Once per scene, when a Wild or Wyrd complication lands on an adjacent ally, draw it onto yourself instead.
    \item[IV.] \textbf{The Still Spot} (once per session, free): for one full round, supernatural effects of the Wild and the Wyrd are suppressed within 3 tiles of you. Things that live in those planes find you \textit{extremely rude}.
\end{description}

\medskip\hrule\medskip

\section{Deeds of the Road}

\subsection{Storied}
\textit{Earned by having the tale arrive in town three days before you did.}
\begin{description}
    \item[I.] +1 die on social rolls with anyone who has heard of you (the GM decides who has; your Reputation line decides what they heard).
    \item[II.] \textbf{An Old Acquaintance} (once per session): declare a contact in any settlement of meaningful size --- a factor, a Rider, a hand you once helped. The GM gives them a name, a debt or a grudge, and a door.
    \item[III.] Your word is credit: once per session, secure goods, passage, or hospitality on Reputation alone, up to the scale of a serious favor. It goes on the Funds line, not the Coins line, and it is \textit{remembered}.
    \item[IV.] \textbf{Precede Yourself} (once per session): when you enter a scene, your legend enters first --- the GM shifts every named NPC's disposition one step (toward awe, fear, or welcome; your Reputation decides which).
\end{description}

\subsection{Steady Hands}
\textit{Earned by fixing, in the dark, under fire, the thing the manual said to replace.}
\begin{description}
    \item[I.] +1 die on Knack.
    \item[II.] \textbf{Jury-Rig} (once per scene): a repair, disable, or rig that would be a long action costs 1 Momentum instead. It holds exactly as long as it needs to.
    \item[III.] Your gear never jams, fouls, or fails at the worst moment, and improvised tools count as proper ones in your hands.
    \item[IV.] \textbf{Exactly the Tool} (once per session): produce from your kit precisely the mundane item this moment requires, up to 5 Rings' worth. You packed it. Obviously you packed it.
\end{description}

\subsection{Cat-Foot}
\textit{Earned by being reported dead, absent, or never there at all --- while present.}
\begin{description}
    \item[I.] +1 die on Nimbleness rolls to move unseen or unheard.
    \item[II.] Moving quietly no longer slows you, and you may hide as part of a Move action.
    \item[III.] \textbf{Vanish} (1 Momentum, once per scene): if any cover or shadow exists, you are unseen; enemies must beat your Nimbleness with Perceive to target you.
    \item[IV.] \textbf{Never There} (once per session): when you are revealed, targeted, or cornered, you were somewhere else all along --- reposition up to your Move in tiles, unseen, as the scene reorganizes around the fact.
\end{description}

\subsection{Long Haul}
\textit{Earned by arriving --- late season, bad road, wrong weather --- anyway.}
\begin{description}
    \item[I.] +1 die against exhaustion, forced march, and privation, and you may travel 2 hours longer per day. (Humans: yes, it stacks. Nobody outlasts you.)
    \item[II.] You can carry a second load --- a companion's pack, a companion --- without penalty.
    \item[III.] Your whole party treats travel conditions as one step gentler while you set the pace.
    \item[IV.] \textbf{In Time} (once per session): retroactively, you pushed through --- the night, the pass, the flood --- and you are where the story needs you at the moment it needs you there, winded and present.
\end{description}

\medskip\hrule\medskip

\section{Table Notes}

\textbf{Six slots is a life.} A veteran holds at most six deeds, and the choice of which legends to deepen versus which to begin is the advancement game: Deadeye IV is a different person from Deadeye II / Quickdraw II, and both are different from the six-at-rank-I generalist whom the frontier calls, with respect, \textit{hard to plan around}.

\textbf{Award ranks out loud.} The best moment to grant a deed rank is thirty seconds after the table saw it earned --- ``write it down: that's Bulwark II'' --- because a deed awarded at the table teaches everyone what the game celebrates. The even-level track exists so no one falls behind; the awards are where the system lives.

\textbf{Name the moment on the line.} The deed's name is shared; its story is not. Players are encouraged to note the earning beside the rank --- \textit{Deadeye III, the Spur, 311} --- because in this world, the Gorûn and the Ostrogar and the Book of the Seen are all correct about the same thing: what you have done is a stat, and it should be written where you were.
